\section{Background and Motivation}
\label{sec:motivation}

\subsection{Workspaces as shared state and private changes}

A workspace contains more than a version-controlled source tree. Tools may
also inspect dependency directories, generated artifacts, build trees, and
temporary files. The intended deployment chooses which directory subtrees
are stable enough to publish as immutable baselines; other state remains
mutable. Multiple tasks can reference the same published subtree without
requiring their subsequent writes to become mutually visible.

The useful workload properties must be established empirically. We will
measure the overlap of files and metadata read across tasks, the fraction of
the initial namespace that each task modifies, and the lifetime of a shared
baseline. These quantities distinguish a favorable workload from one in which
tasks rewrite most files or repeatedly invalidate their dependencies. Agent
framework papers establish that repository navigation and editing occur;
they do not establish the frequency of \sysname's storage bottlenecks.
We therefore treat high sharing and sparse mutation as hypotheses rather
than facts about all coding agents.

An earlier project pilot used DSH 0.1.5-rc.1 with the configured model
\code{deepseek-v4-flash} over a WSL passthrough FUSE file system. It is a
workload lead, not a \sysname/RustFS performance measurement. This draft does
not reuse its numerical results because the raw event records and timing
definitions have not been revalidated here. The formal study will recover or
recapture the traces and distinguish handler service time, overlapping I/O,
tool time, and task wall-clock time.

\subsection{Metadata sharing does not determine metadata access}

Namespace organization and caching have long been central to file-system
performance. TableFS organizes local metadata in a table;
IndexFS combines partitioning, caching, and bulk insertion; LocoFS and
InfiniFS revisit the organization of distributed metadata
services~\cite{tablefs2013,indexfs2014,locofs2017,infinifs2022}.
These systems caution against assuming that every metadata operation always
requires a distinct network round trip. A fair comparison must include their
relevant techniques and an equal cache budget.

\sys studies a restricted opportunity: selected lower subtrees cannot change
after publication. Their lookup records need not participate in the same
update protocol as the private upper. An immutable representation can carry
directory entries, frequently used attributes, and data locators together,
allowing a client to fetch enough information for several related operations.
The benefit depends on access locality and cache state. For random point
lookups, a group may fetch more bytes than an individual record. For a cold
deep path, several index pages may be required. Both cases belong in the
motivation study and the evaluation.

The property being tested is \emph{lower-query independence}, not the
elimination of a metadata service. Upper lookups, workspace control, reader
pins, publication, and reclamation may still require mutable records.
This distinction is also essential when comparing \sys with DeltaFS, whose
private logs and selectively assembled namespaces already challenge a
single global metadata ground truth~\cite{deltafs2021}.

\subsection{Logical request size differs from physical work}

There are at least five relevant sizes: an application's request, a
file-system logical chunk, the finalized uploaded object, an erasure-coding
block, and a per-shard verification/read frame. Compression, framing, and
partial tail objects make the uploaded length differ from file length.
Erasure coding then distributes each coding block across data and parity
shards. A range request can require a larger verified interval on each shard
that serves it.

Storage systems already make different decisions about striping, coding,
and redundancy transitions~\cite{azureec2012,morph2024,okapi2025}.
\sysname's proposed connection is specifically between finalized object
length and the backend's coding and verification geometry. It leaves the
data/parity ratio fixed during a causal comparison. Smaller blocks do not
by themselves reduce that ratio, and smaller verification frames add
checksums and processing overhead.

Moreover, the client may fetch a full object even for a short application
read. In that case, a smaller backend quantum may not reduce the bytes
requested. We will trace the complete path, including the lower reader,
upper integrity check, client cache, HTTP request, shard reads, and repair.
This is necessary to determine whether a physical-layout optimization
benefits the intended workspace workload.
